\documentclass[10pt,a4paper]{amsart}
\usepackage[margin=19mm]{geometry}
\usepackage{amsmath,amssymb,mathtools}
\usepackage[scr=rsfs,cal=euler]{mathalfa}
\usepackage{xfrac}
\usepackage[unicode,hidelinks,bookmarks=false]{hyperref}
\newcommand{\ie}{i.e.\ }
\newcommand{\cf}{cf.\ }
\newcommand{\resp}{respectively\ }
\newcommand{\Subsection}{\subsection}
\providecommand{\nobreakdash}{\nobreak\hbox{-}}
\setlength{\emergencystretch}{2em}
\multlinegap0pt
\usepackage{amssymb}
\usepackage{dsfont}
\newtheorem{assumption}{Assumption}
\newtheorem{theorem}{Theorem}
\newtheorem{proposition}{Proposition}
\def\EE{{\mathcal E}}
\newcommand{\Ee}{\mathds E}
\newcommand{\F}{\mathscr{F}}
\def\FF{{\mathcal F}}
\newcommand{\HH}{\mathcal{H}}
\newcommand{\Pp}{\mathds P}
\newcommand{\R}{\mathds R}
\newcommand{\e}{\varepsilon}
\renewcommand{\ge}{\geqslant}
\renewcommand{\le}{\leqslant}
\renewcommand{\leq}{\leqslant}
\newcommand{\w}{\omega}
\title{Stochastic homogenization of stable-like process with divergence free drift}
\author{Xin Chen, Kun Yin}
\date{Source: arXiv:2505.15300v2 (CC BY 4.0)}
\begin{document}
\maketitle

\noindent\emph{Source and licence.} Stochastic homogenization of stable-like process with divergence free drift. Version arXiv:2505.15300v2 (CC BY 4.0), distributed by arXiv under the Creative Commons Attribution 4.0 International licence, http://creativecommons.org/licenses/by/4.0/, as recorded in the arXiv licence metadata for this exact version and shown on the abstract page https://arxiv.org/abs/2505.15300v2. Full licence notice: \texttt{licenses/arxiv-2505.15300-CC-BY-4.0.txt}.\par\smallskip

\noindent\textit{Editorial scope.} Counted unit: the article's theorem t1-2 (source0.tex lines 449--461). The article's complete original proof of it is delivered with it (source0.tex lines 697--1007, one \texttt{proof} environment of 311 lines, which the article places away from the statement). No mathematics is written, repaired or re-derived: the statement and the proof are the article's own lines, in the article's own notation, and no retained line is substituted or re-flowed.  Retained uncounted as the source-local context the proof needs: lines 392--405; lines 441--443; lines 471--490.  Nothing was omitted from any retained span, so the package carries no omission record.  The retained lines cite 2 works, whose entries are transcribed verbatim from the article's own bibliography source in the e-print and delivered with short editorial labels recorded in the item's reference mapping.  No retained line contains a figure environment, an image-inclusion command or drawing code, so no image is delivered and the package declares no asset dependency.  Editorial formatting only, in addition: an AMS \texttt{amsart} preamble that restates the article's own declarations and macros used by the retained lines, namely the environments \texttt{assumption}, \texttt{theorem}, \texttt{proposition} on separate counters, together with the standard packages those lines need.  The retained environments print in this document as Assumption 1, Theorem 1, Proposition 1 in order of appearance, whereas the article prints the same items with the numbering of its own full document.  The complete native source of the article as distributed in the arXiv e-print for this exact version is bundled unchanged as \texttt{sources/178500/source0.tex}.
\begin{assumption}\label{a1-2}
$\tilde b=(\tilde b_1,\tilde b_2,\cdots, \tilde b_d):\Omega \to \R^d$ is a random vector
and there exists a system of stream functions(for $\tilde b$) $\tilde H_{jl}:\Omega \to \R$, $1\le j,l\le d$ such that
the following properties hold.
\begin{itemize}
\item [(i)] $\tilde H_{jl}\in \HH_{1}$ for every $1\le j,l \le d$;

\item [(ii)] $\tilde H_{jl}=-\tilde H_{lj}$ for every $1\le j,l \le d$;

\item [(iii)] $\tilde b_j=\sum_{l=1}^d D_l \tilde H_{jl}$ for every $1\le j \le d$;

\item [(iv)] The function $x\mapsto D_l \tilde H_{jl}(\tau_x \w)$ is continuous for every $1\le l \le d$ and a.s. $\w\in \Omega$.
\end{itemize}
\end{assumption}

\begin{equation}\label{e1-6}
\EE^{\e,\w} (u_{\lambda,f}^{\e,\w},g)+\lambda\int_{\R^d}u_{\lambda,f}^{\e,\w}(x)g(x)dx=\int_{\R^d}f(x)g(x)dx,\ \forall\ g\in C_c^1(\R^d).
\end{equation}

\begin{proposition}\label{t1-1}
For every $f\in C_c^1(\R^d)$, $\lambda>0$, $\e\in (0,1)$ and a.s. $\w\in \Omega$, there exists a $u_{\lambda,f}^{\e,\w}\in \FF_0^{\e,\w}$ such that
\eqref{e1-6} holds,
%\begin{equation}\label{t1-1-1}
%\EE^{\e,\w} (u_{\lambda,f}^{\e,\w},g)+\lambda\int_{\R^d}u_{\lambda,f}^{\e,\w}(x)g(x)dx=\int_{\R^d}f(x)g(x)dx,\ \forall\ g\in C_c^1(\R^d),
%\end{equation}
and

\begin{equation}\label{t1-1-1a}
\sup_{\e\in (0,1)}\|u_{\lambda,f}^{\e,\w}\|_{L^2(\R^d)}<+\infty,\ \sup_{\e\in (0,1)}\|u_{\lambda,f}^{\e,\w}\|_{W^{\alpha/2,2}(\R^d)}<+\infty,\ \
\sup_{\e\in (0,1)}\EE_0^{\e,\w}\left(u_{\lambda,f}^{\e,\w},u_{\lambda,f}^{\e,\w}\right)<+\infty.
\end{equation}

Moreover if we assume $\alpha\in [1,2)$ and

\begin{equation}\label{t1-1-2a}
\sup_{\w\in \Omega}\left(|\tilde b(\w)|+\int_{\R^d}\frac{|b(\tau_z \w)-b(\w)|^2}{|z|^{d+2-\alpha}}dz\right)<+\infty,
\end{equation}
then $u_{\lambda,f}^{\e,\w}$ above is the unique solution of \eqref{e1-6} in $\FF_0^{\e,\w}$.
\end{proposition}

\begin{theorem}\label{t1-2}
For every $f\in C_c^1(\R^d)$, $\lambda>0$ and a.s. $\w\in \Omega$, let $u_{\lambda,f}^{\e,\w}$ be the solution
of \eqref{e1-6} obtained in Proposition \ref{t1-1}, then we have
\begin{equation}\label{t1-2-1}
\lim_{\e \to 0}\int_{O}|u_{\lambda,f}^{\e,\w}(x)-\bar u_{\lambda,f}(x)|^2dx=0,\ {\rm for\ every\ bounded\ subset\ } O\subset \R^d.
\end{equation}
Here $\bar u_{\lambda,f}$ is the unique solution in $L^2(\R^d)$ of the following equation
\begin{equation}\label{t1-2-2}
\lambda \bar u_{\lambda,f}(x)-\bar L \bar u_{\lambda,f}(x)=f(x),\ \ x\in \R^d,
\end{equation}
with
$$\bar L f(x):={\rm p.v.}\int_{\R^d}\left(f(x+z)-f(x)\right)\frac{\Ee[\mu]^2}{|z|^{d+\alpha}}dz,\ \ f\in C_c^2(\R^d).$$
\end{theorem}

\begin{proof} [Proof of Theorem $\ref{t1-2}$]
By \eqref{t1-1-1a} we know $\{u_{\lambda,f}^{\e,\w}\}_{\e\in (0,1)}$ is weakly compact in
$W^{\alpha/2,2}(\R^d)$, hence it is compact in $L_{loc}^2(\R^d)$. Then it suffices to prove that for every
subsequence $\{u_{\lambda,f}^{\e_m,\w}\}_{m\ge 1}$ which is convergent in $L_{loc}^2(\R^d)$, it holds that
\begin{equation*}
\lim_{m \to \infty}\int_{\R^d}\left|u_{\lambda,f}^{\e_m,\w}-\bar u_{\lambda,f}(x)\right|^2dx=0, \ \ {\rm for\ every\ bounded\ subset}\ O\subset \R^d,
\end{equation*}
where $\bar u_{\lambda,f}$ is the unique solution of \eqref{t1-2-2}.

Suppose that $\{u_{\lambda,f}^{\e_m,\w}\}_{m\ge 1}$ is a (arbitrarily fixed) subsequence such that
\begin{equation}\label{t1-2-3}
\lim_{m \to \infty}\int_{O}\left|u_{\lambda,f}^{\e_m,\w}(x)-u_0^{\w}(x)\right|^2dx=0,\ \ {\rm for\ every\ bounded\ subset}\ O\subset \R^d,
\end{equation}
for some $u_0^{\w}\in L^2(\R^d)$. Now we are going to prove $u_0^{\w}=\bar u_{\lambda,f}$.

According to \eqref{e1-6} we obtain for every $g\in C_c^1(\R^d)$,
\begin{equation}\label{t1-2-4}
\EE_0^{\e_m,\w}(u_{\lambda,f}^{\e_m,\w},g)-
\e_m^{-(\alpha-1)}\sum_{j=1}^d\int_{\R^d}b_j\left(\frac{x}{\e},\w\right)\frac{\partial g(x)}{\partial x_j}u_{\lambda,f}^{\e_m,\w}(x)dx
+\lambda\int_{\R^d}u_{\lambda,f}^{\e_m,\w}(x)g(x)dx=\int_{\R^d}f(x)g(x)dx.
\end{equation}

{\bf Step 1} In this step, we are going to prove
\begin{equation}\label{t1-2-5}
\lim_{m \to \infty}\EE_0^{\e_m,\w}(u_{\lambda,f}^{\e_m,\w},g)=-\int_{\R^d}u_0^\w(x)\bar L g(x)dx,\ \forall\ g\in C_c^2(\R^d).
\end{equation}
The main method in this step is similar with that in \cite{CCKW1}, for convenience of the reader we will give a detailed proof here.
For every $\delta\in (0,1)$ we set
\begin{equation*}
\begin{split}
&\quad \EE_0^{\e_m,\w}(u_{\lambda,f}^{\e_m,\w},g)\\
&=\frac{1}{2}
\int_{\R^d}\left(\int_{\{|z|\le \delta\}}+\int_{\{\delta<|z|<\delta^{-1}\}}+\int_{\{|z|\ge \delta^{-1}\}}\right)
\left(u_{\lambda,f}^{\e_m,\w}(x+z)-u_{\lambda,f}^{\e_m,\w}(x)\right)\left(g(x+z)-g(x)\right)\\
&\cdot\frac{\mu\left(\frac{x}{\e_m};\w\right)\mu\left(\frac{x+z}{\e_m};\w\right)}{|z|^{d+\alpha}}dzdx
=:I_1^{m,\delta,\w}+I_2^{m,\delta,\w}+I_3^{m,\delta,\w}.
\end{split}
\end{equation*}
Suppose that ${\rm supp}g\subset B(0,R_0):=\{x\in \R^d; |x|\le R_0\}$ for some $R_0>0$. Apply Cauchy-Schwartz inequality
we derive that
\begin{align*}
|I_1^{m,\delta,\w}|&\le \left(\frac{1}{2}\int_{\R^d}\int_{\{|z|\le \delta\}}
\left(u_{\lambda,f}^{\e_m,\w}(x+z)-u_{\lambda,f}^{\e_m,\w}(x)\right)^2\frac{\mu\left(\frac{x}{\e_m};\w\right)\mu\left(\frac{x+z}{\e_m};\w\right)}{|z|^{d+\alpha}}dzdx\right)^{1/2}\\
&\quad \cdot \left(\frac{1}{2}\int_{\R^d}\int_{\{|z|\le \delta\}}
\left(g(x+z)-g(x)\right)^2\frac{\mu\left(\frac{x}{\e_m};\w\right)\mu\left(\frac{x+z}{\e_m};\w\right)}{|z|^{d+\alpha}}\right)^{1/2}\\
&\le \EE_0^{\e_m,\w}\left(u_{\lambda,f}^{\e_m,\w},u_{\lambda,f}^{\e_m,\w}\right)^{1/2}\cdot
\left(\frac{\|\nabla g\|_\infty^2}{2}\int_{\{|x|\le R_0+1\}}\int_{\{|z|\le \delta\}}
\left(\mu\left(\frac{x+z}{\e_m};\w\right)^2+\mu\left(\frac{x}{\e_m};\w\right)^2\right)\frac{|z|^2}{|z|^{d+\alpha}}dzdx\right)^{1/2}\\
&\le c_1(\w)
\left(\|\nabla g\|_\infty^2\int_{\{|x|\le R_0+2\}}\mu\left(\frac{x}{\e_m};\w\right)^2\left(\int_{\{|z|\le \delta\}}
\frac{|z|^2}{|z|^{d+\alpha}}dz\right)dx\right)^{1/2},
\end{align*}
where $c_1(\w):=\sup_{m\ge 1}\EE_0^{\e_m,\w}\left(u_{\lambda,f}^{\e_m,\w},u_{\lambda,f}^{\e_m,\w}\right)^{1/2}$
is a positive constant(may depend on $\w$) which is finite due to \eqref{t1-1-1a}.

By direct computation we obtain
\begin{align*}
\int_{\{|x|\le R_0+2\}}\int_{\{|z|\le \delta\}}
\mu\left(\frac{x}{\e_m};\w\right)^2\frac{|z|^2}{|z|^{d+\alpha}}dzdx&\le
c_2\delta^{2-\alpha}\cdot \int_{\{|x|\le R_0+2\}}\tilde \mu\left(\tau_\frac{x}{\e_m}\w\right)^2dx.
\end{align*}
Hence by Birkhoff ergodic theorem (see, for example, \cite[Proposition 2.1]{CCKW1} or \cite[Theorem 7.2]{JKO}) it holds that for every fixed $\delta\in (0,1)$,
\begin{align*}
\lim_{m \to \infty}\int_{\{|x|\le R_0+2\}}\int_{\{|z|\le \delta\}}
\mu\left(\frac{x}{\e_m};\w\right)^2\frac{|z|^2}{|z|^{d+\alpha}}dzdx \le c_3\Ee[\tilde \mu^2]\delta^{2-\alpha}.
\end{align*}



Similarly for $I_3^{m,\delta,\w}$ we obtain
\begin{align*}
|I_3^{m,\delta,\w}|&\le \left(\frac{1}{2}\int_{\R^d}\int_{\{|x-y|> \delta^{-1}\}}
\left(u_{\lambda,f}^{\e_m,\w}(y)-u_{\lambda,f}^{\e_m,\w}(x)\right)^2\frac{\mu\left(\frac{x}{\e_m};\w\right)\mu\left(\frac{y}{\e_m};\w\right)}{|x-y|^{d+\alpha}}dydx\right)^{1/2}\\
&\quad \cdot \left(\frac{1}{2}\int_{\R^d}\int_{\{|x-y|> \delta^{-1}\}}
\left(g(y)-g(x)\right)^2\frac{\mu\left(\frac{x}{\e_m};\w\right)\mu\left(\frac{y}{\e_m};\w\right)}{|x-y|^{d+\alpha}}dydx\right)^{1/2}.\\
%&\le \EE_0^{\e_m,\w}\left(u_{\lambda,f}^{\e_m,\w},u_{\lambda,f}^{\e_m,\w}\right)^{1/2}\cdot
%\left(\|g\|_\infty^2\int_{\{|x|\le R_0+1\}}\int_{\{|z|> \delta^{-1}\}}
%\frac{\left(\mu\left(\frac{x+z}{\e_m};\w\right)^2+\mu\left(\frac{x}{\e_m};\w\right)^2\right)}{|z|^{d+\alpha}}dzdx\right)^{1/2}\\
\end{align*}
For $\delta^{-1}>2(R_0+2)$, we know $\left(g(x)-g(y)\right)1_{\{|x-y|>\delta^{-1}\}}\neq 0$ only if
$x\in B(0,R_0)$ or $y\in B(0,R_0)$. Then using the symmetry of $x$ and $y$ we have
\begin{align*}
&\varlimsup_{m\to \infty} \iint_{\{|x-y|>\delta^{-1}\}}(g(y)- g(x))^2\frac{\mu\left(\frac{x}{\e_m};\w\right)\mu\left(\frac{y}{\e_m};\w\right)}{|x-y|^{d+\alpha}}\,dx\,dy
\\
%&\leq  2 \limsup_{\e \to 0}  \iint_{\{|x-y|>1/\eta\}}(g(y)- g(x))^2\frac{\mu\left(\frac{x}{\e_m};\w\right)\mu\left(\frac{y}{\e_m};\w\right)}{|x-y|^{d+\alpha}}\,dx\,dy \\
& \leq  4\|g\|_\infty^2\varlimsup_{m \to \infty}  \int_{B(0, R_0)} \left( \int_{\{|y-x|>\delta^{-1}\}} \frac{\tilde \mu\left(\frac{y}{\e_m};\w\right)}
{|y-x|^{d+\alpha}} dy \right)
 \mu\left(\frac{x}{\e_m};\w\right) \,dx \\
&\le c_2\varlimsup_{m \to \infty}\int_{\{|x|\le R_0+2\}}\mu\left(\frac{x}{\e_m};\w\right)^2\left(\int_{\{|x-y|> \delta^{-1}\}}
\frac{1}{|x-y|^{d+\alpha}}dy\right)dx\\
&\quad\quad +c_2\varlimsup_{m \to \infty}\int_{\{|x|\le R_0+2\}}\int_{\{|z|> \delta^{-1}\}}\frac{\mu\left(\frac{y}{\e_m};\w\right)^2}{|x-y|^{d+\alpha}}dydx,
\end{align*}
where the last step follows from Cauchy-Schwartz inequality.
Applying Birkhoff ergodic theorem  we get
\begin{align*}
\lim_{m \to \infty}\int_{\{|x|\le R_0+2\}}\mu\left(\frac{x}{\e_m};\w\right)^2\left(\int_{\{|x-y|> \delta^{-1}\}}
\frac{1}{|x-y|^{d+\alpha}}dy\right)dx\le c_4\Ee[\tilde \mu^2]\delta^{\alpha}.
\end{align*}
%By change of variable $y=x+z$,
For $\delta^{-1}\ge 2(R_0+2)$ we have
\begin{align*}
%\quad\int_{\{|x|\le R_0+2\}}\int_{\{|z|> \delta^{-1}\}}\frac{\mu\left(\frac{x+z}{\e_m};\w\right)^2}{|z|^{d+\alpha}}dzdx
\int_{\{|x|\le R_0+2\}}\int_{\{|x-y|> \delta^{-1}\}}\frac{\tilde \mu\left(\tau_{\frac{y}{\e_m}}\w\right)^2}{|x-y|^{d+\alpha}}dydx
&\le c_5\int_{\{|x|\le R_0+2\}}\int_{\{|y|> (2\delta)^{-1}\}}\frac{\tilde \mu\left(\tau_{\frac{y}{\e_m}}\w\right)^2}{|y|^{d+\alpha}}dydx\\
&=:c_5\sum_{k=0}^\infty |B(0,R_0+2)|\int_{U_k}\frac{\tilde \mu\left(\tau_{\frac{y}{\e_m}}\w\right)^2}{|y|^{d+\alpha}}dy\\
&\le c_6\delta^{d+\alpha}\sum_{k=0}^\infty 2^{-k(d+\alpha)}\e_m^{d}\int_{B\left(0,\frac{2^{k+1}(2\delta)^{-1}}{\e_m}\right)}\tilde \mu\left(\tau_{y}\w\right)^2 dy.
\end{align*}
where $U_k:=\{y\in \R^d; 2^k(2\delta)^{-1}<|y|\le 2^{k+1}(2\delta)^{-1}\}$ for every integer $k\ge 0$. According to the ergodic theorem again
there exists a constant $c_7(\w)$(which may depend on $\w$) such that
\begin{align*}
\sup_{m\ge 1}\e_m^d\int_{B(0,\e_m^{-1})} \tilde \mu(\tau_y \w)^2 dy\le c_7(\w)\Ee[\tilde \mu^2].
\end{align*}
Hence using this in above inequality yields
\begin{align*}
\sup_{m\ge 1}\int_{\{|x|\le R_0+2\}}\int_{\{|x-y|> \delta^{-1}\}}\frac{\mu\left(\frac{y}{\e_m};\w\right)^2}{|x-y|^{d+\alpha}}dydx
&\le c_8(\w)\delta^\alpha\Ee[\tilde \mu^2]\sum_{k=0}^\infty 2^{-k\alpha}\le c_9(\w)\delta^\alpha.
\end{align*}
Combining all above estimates we can verify directly that
\begin{align}\label{t1-2-6}
\lim_{\delta \downarrow 0}\varlimsup_{m \to \infty}\left(|I_1^{m,\delta,\w}|+|I_3^{m,\delta,\w}|\right)=0.
\end{align}




By the change of variable $\tilde z=-z$ and $\tilde x=x+z$ for the integral
$$\int_{\R^d}\int_{\{\delta<|z|<\delta^{-1}\}}u_{\lambda,f}^{\e_m,\w}\left(x+z;\w\right)\left(g(x+z)-g(x)\right)
\frac{\mu\left(\frac{x}{\e_m};\w\right)\mu\left(\frac{x+z}{\e_m};\w\right)}{|z|^{d+\alpha}}dzdx$$
we have
\begin{align*}
I_2^{m,\delta,\w}&=-\int_{\R^d}\int_{\{\delta<|z|<\delta^{-1}\}}
u_{\lambda,f}^{\e_m,\w}\left(x;\w\right)\left(g(x+z)-g(x)\right)
\frac{\mu\left(\frac{x}{\e_m};\w\right)\mu\left(\frac{x+z}{\e_m};\w\right)}{|z|^{d+\alpha}}dzdx,
\end{align*}
which implies that for each fixed $\delta\in (0,1)$,
\begin{align*}
&\quad \lim_{m \to \infty}\left|I_2^{m,\delta,\w}-\frac{1}{2}\int_{\R^d}\int_{\{\delta<|z|<\delta^{-1}\}}
\left(u_{0}^{\w}(x+z)-u_0^\w(x)\right)\left(g(x+z)-g(x)\right)
\frac{\mu\left(\frac{x}{\e_m};\w\right)\mu\left(\frac{x+z}{\e_m};\w\right)}{|z|^{d+\alpha}}dzdx\right|\\
&\le \lim_{m \to \infty}c_{10}\|g\|_\infty\delta^{-d-\alpha}\cdot
\left(\int_{\{|x|\le R_0+\delta^{-1}\}}\int_{\{\delta<|z|<\delta^{-1}\}}\left|u_{\lambda,f}^{\e_m,\w}\left(x\right)-u_0^\w(x)\right|^2dzdx\right)^{1/2}\\
&\cdot \left(\int_{\{|x|\le R_0+\delta^{-1}\}}\int_{\{\delta<|z|<\delta^{-1}\}}\mu\left(\frac{x}{\e_m};\w\right)^2
\mu\left(\frac{x+z}{\e_m};\w\right)^2 dzdx\right)^{1/2}\\
&\le \lim_{m \to \infty}c_{11}(\delta)
\left(\int_{\{|x|\le R_0+\delta^{-1}\}}\left|u_{\lambda,f}^{\e_m,\w}\left(x\right)-u_0^\w(x)\right|^2dx\right)^{1/2}\\
&\cdot \left(\int_{\{|x|\le R_0+2\delta^{-1}\}}\int_{\{|y|\le R_0+2\delta^{-1}\}}\tilde \mu\left(\tau_{\frac{x}{\e_m}}\w\right)^2
\tilde \mu\left(\tau_{\frac{y}{\e_m}}\w\right)^2 dxdy\right)^{1/2}\\
&\le \lim_{m \to \infty}c_{11}(\delta)\Ee[\tilde \mu^2]\left(\int_{\{|x|\le R_0+\delta^{-1}\}}\left|u_{\lambda,f}^{\e_m,\w}\left(x\right)-u_0^\w(x)\right|^2dx\right)^{1/2}=0,
\end{align*}
where in the third inequality we have used the ergodic theorem, and the last step is due to \eqref{t1-2-3}.

According to \cite[Lemma 3.1(ii)]{CCKW1} we have
\begin{align*}
&\quad \lim_{m \to \infty}\int_{\R^d}\int_{\{\delta<|z|<\delta^{-1}\}}
\left(u_{0}^{\w}(x+z)-u_0^\w(x)\right)\left(g(x+z)-g(x)\right)
\frac{\mu\left(\frac{x}{\e_m};\w\right)\mu\left(\frac{x+z}{\e_m};\w\right)}{|z|^{d+\alpha}}dzdx\\
&=\int_{\R^d}\int_{\{\delta<|z|<\delta^{-1}\}}
\left(u_{0}^{\w}(x+z)-u_0^\w(x)\right)\left(g(x+z)-g(x)\right)
\frac{\Ee[\tilde \mu]^2}{|z|^{d+\alpha}}dzdx.
\end{align*}
Therefore putting all these estimates together yields that
\begin{align*}
\lim_{m \to \infty}I_2^{m,\delta,\w}&=\frac{1}{2}\int_{\R^d}\int_{\{\delta<|z|<\delta^{-1}\}}
\left(u_{0}^{\w}(x+z)-u_0^\w(x)\right)\left(g(x+z)-g(x)\right)
\frac{\Ee[\tilde \mu]^2}{|z|^{d+\alpha}}dzdx\\
&=-\int_{\R^d}u_0^{\w}(x)\left(\int_{\{\delta<|z|<\delta^{-1}\}}\left(g(x+z)-g(x)\right)
\frac{\Ee[\tilde \mu]^2}{|z|^{d+\alpha}}dz\right)dx.
\end{align*}
According to this and \eqref{t1-2-6}, firstly letting $m \to \infty$, then $\delta \downarrow 0$ we
can prove \eqref{t1-2-5}.

{\bf Step 2} It remains to prove that
\begin{align}\label{t1-2-7}
\lim_{m \to \infty}\e_m^{-(\alpha-1)}\sum_{j=1}^d\int_{\R^d}b_j\left(\frac{x}{\e_m};\w\right)\frac{\partial g(x)}{\partial x_j}u_{\lambda,f}^{\e_m,\w}(x)dx=0,
\ \forall\ g\in C_c^2(\R^d).
\end{align}
Without loss of generality, we assume that ${\rm supp}g\subset B(0,R_0)$ for some $R_0>0$.

%When $\alpha\in (0,1)$, according to the fact $\tilde b_j \in L^2(\Omega;\Pp)$ and \eqref{t1-1-1a} and applying Cauchy-Schwartz inequality and
%ergodic theorem we obtain
%\begin{align*}
%&\quad\lim_{m \to \infty}\e_m^{-(\alpha-1)}\sum_{j=1}^d\left|\int_{\R^d}b_j\left(\frac{x}{\e_m},\w\right)\frac{\partial g(x)}{\partial x_j}u_{\lambda,f}^{\e_m,\w}(x)dx\right|\\
%&\le c_{12}\sup_{\e>0}\|u_{\lambda,f}^{\e,\w}\|_{L^2(\R^d)}\cdot\lim_{m \to \infty}\e_m^{1-\alpha}
%\sqrt{\int_{B(0,R_0)}\left|\tilde b_j\left(\tau_{\frac{x}{\e_m}}\w\right)\right|^2dx}=0.
%\end{align*}
%This means that \eqref{t1-2-7} is true for $\alpha\in (0,1)$.

%Now we consider the case $\alpha\in [1,2)$.
By Assumption \ref{a1-2} we have $H_{jl}(\cdot;\w)\in W^{1,2}_{loc}(\R^d)$ for $H_{jl}(x;\w):=\tilde H_{jl}(\tau_x \w)$.
%and $b_j\left(x;\w\right)=\sum_{l=1}^d \frac{\partial H_{jl}(\cdot;\w)}{\partial x_l}(x)$.
Since ${\rm supp}g\subset B(0,R_0)$, we choose a function $h\in C_c^2(\R^d)$ such that
$h(x)=1$ for every $x\in B(0,R_0)$ and ${\rm supp}h \subset B(0,2R_0)$. Based on this it holds that
\begin{align}\label{t1-2-7a}
\e_m^{-(\alpha-1)}\int_{\R^d}b_j\left(\frac{x}{\e},\w\right)\frac{\partial g(x)}{\partial x_j}u_{\lambda,f}^{\e_m,\w}(x)dx&=
\e_m^{2-\alpha}\sum_{l=1}^d\int_{\R^d}  \frac{\partial}{\partial x_l}\left( H_{jl}^{\e_m}\left(\cdot;\w\right)h(\cdot)\right)(x)
\frac{\partial g(x)}{\partial x_j}u_{\lambda,f}^{\e_m,\w}(x)dx,
\end{align}
where $H_{jl}^{\e_m}(x;\w):=H_{jl}\left(\frac{x}{\e_m};\w\right)$.

So by Parseval's identity we get for every $1\le j,l\le d$,
\begin{align*}
&\quad \e_m^{2-\alpha}\left|\int_{\R^d}\frac{\partial}{\partial x_l}\left( H_{jl}^{\e_m}\left(\cdot;\w\right)h(\cdot)\right)(x)
\frac{\partial g(x)}{\partial x_j}u_{\lambda,f}^{\e_m,\w}(x)dx\right|\\
&=\e_m^{2-\alpha}\left|\int_{\R^d}i \xi_l \F\left(H_{jl}^{\e_m}\left(\cdot;\w\right)h(\cdot)\right)(\xi)
\cdot\F\left(\frac{\partial g(\cdot)}{\partial x_j}u_{\lambda,f}^{\e_m,\w}(\cdot)\right)(\xi)d\xi\right|\\
&\le c_{13}\e_m^{2-\alpha}\left(\int_{\R^d}|\xi|^\alpha\left|\F\left(\frac{\partial g(\cdot)}{\partial x_j}
u_{\lambda,f}^{\e_m,\w}(\cdot)\right)(\xi)\right|^2d\xi\right)^{1/2}\cdot
\left(\int_{\R^d}|\xi|^{2-\alpha}\left|\F\left(H_{jl}^{\e_m}\left(\cdot;\w\right)h(\cdot)\right)(\xi)\right|^2d\xi\right)^{1/2}.
\end{align*}
Then we have
\begin{align*}
&\quad\int_{\R^d}|\xi|^\alpha \left|\F\left(\frac{\partial g(\cdot)}{\partial x_j}
u_{\lambda,f}^{\e_m,\w}(\cdot)\right)(\xi)\right|^2d\xi\\
&\le c_{14}\int_{\R^d}\int_{\R^d}\frac{\left|\frac{\partial g(x+z)}{\partial x_j}
u_{\lambda,f}^{\e_m,\w}(x+z)-\frac{\partial g(x)}{\partial x_j}
u_{\lambda,f}^{\e_m,\w}(x)\right|^2}{|z|^{d+\alpha}}dzdx\\
&\le c_{15}\int_{\R^d}\int_{\R^d}\frac{|u_{\lambda,f}^{\e_m,\w}(x+z)-u_{\lambda,f}^{\e_m,\w}(x)|^2}{|z|^{d+\alpha}}
\left|\frac{\partial g(x+z)}{\partial x_j}\right|^2dzdx
+c_{15}\int_{\R^d}|u_{\lambda,f}^{\e_m,\w}(x)|^2\left(\int_{\R^d}\frac{\left|\frac{\partial g(x+z)}{\partial x_j}-\frac{\partial g(x)}{\partial x_j}\right|^2}{|z|^{d+\alpha}}dz\right)dx\\
&\le c_{16}\int_{\R^d}\int_{\R^d}\frac{|u_{\lambda,f}^{\e_m,\w}(x+z)-u_{\lambda,f}^{\e_m,\w}(x)|^2}{|z|^{d+\alpha}}
dzdx+
%c_{15}\int_{\R^d}\frac{|u_{\lambda,f}^{\e_m,\w}(x)|^2}{(1+|x|)^{d+\alpha}}dx.
c_{15}\int_{\R^d}|u_{\lambda,f}^{\e_m,\w}(x)|^2dx.
\end{align*}
Hence according to \eqref{t1-1-1a} we obtain
\begin{equation}\label{t1-2-9}
\sup_{m\ge 1}\int_{\R^d}|\xi|^\alpha \left|\F\left(\frac{\partial g(\cdot)}{\partial x_j}
u_{\lambda,f}^{\e_m,\w}(\cdot)\right)(\xi)\right|^2d\xi<+\infty.
\end{equation}

Meanwhile it holds that
\begin{align*}
&\quad\int_{\R^d}|\xi|^{2-\alpha}\left|\F\left(H_{jl}^{\e_m}(\cdot;\w)h(\cdot)\right)(\xi)\right|^2d\xi\\
&=\int_{\R^d}\int_{\R^d}\frac{\left|H_{jl}\left(\frac{x}{\e_m};\w\right)h(x)-H_{jl}\left(\frac{y}{\e_m};\w\right)h(y)\right|^2}{|x-y|^{d+2-\alpha}}dxdy\\
&\le 2\int_{\R^d}\int_{\R^d}\frac{\left|H_{jl}\left(\frac{x}{\e_m};\w\right)-H_{jl}\left(\frac{y}{\e_m};\w\right)\right|^2|h(y)|^2}{|x-y|^{d+2-\alpha}}dxdy
+2\int_{\R^d}\int_{\R^d}\frac{|h(x)-h(y)|^2}{|x-y|^{d+2-\alpha}}\left|H_{jl}\left(\frac{x}{\e_m};\w\right)\right|^2dxdy\\
&\le c_{17}\int_{\R^d}\int_{\R^d}\frac{\left|H_{jl}\left(\frac{x}{\e_m};\w\right)-H_{jl}\left(\frac{y}{\e_m};\w\right)\right|^2|h(y)|^2}{|x-y|^{d+2-\alpha}}dxdy
+c_{17}\int_{\R^d}\frac{\left|H_{jl}\left(\frac{x}{\e_m};\w\right)\right|^2}{(1+|x|)^{d+2-\alpha}}dx,
\end{align*}
where the last step we have used the property(since $h\in C_c^\infty(\R^d)$)
\begin{align*}
\quad \int_{\R^d}\frac{\left|h(x)-h(y)\right|^2}{|x-y|^{d+2-\alpha}}dy
&\le \|\nabla h\|_\infty^2\cdot \int_{\{y\in \R^d; |y-x|\le 1\}}\frac{|y-x|^2}{|y-x|^{d+2-\alpha}}dz\\
&+4\|h\|_\infty^2 \cdot \int_{\{y\in \R^d;|y-x|>1\}}\frac{1}{|x-y|^{d+2-\alpha}}dz\le c_{16},\ \forall x\in B(0,3R_0),
\end{align*}
\begin{align*}
\int_{\R^d}\frac{\left|h(x)-h(y)\right|^2}{|x-y|^{d+2-\alpha}}dx
&\le \|h\|_\infty^2\cdot\left(\int_{\{y\in \R^d; |y|\le 2R_0\}}\frac{1}{|x-y|^{d+2-\alpha}}dy\right)\le c_{17}(1+|x|)^{-d-2+\alpha},\ \forall\ x\notin B(0,3R_0).
\end{align*}
By change of variable we obtain
\begin{align*}
&\quad \int_{\R^d}\int_{\R^d}\frac{\left|H_{jl}\left(\frac{x}{\e_m};\w\right)-H_{jl}\left(\frac{y}{\e_m};\w\right)\right|^2|h(y)|^2}{|x-y|^{d+2-\alpha}}dxdy\\
&\le
\e_m^{d-(2-\alpha)}\|h\|_\infty^2\int_{B\left(0,\frac{2R_0}{\e_m}\right)}\left(\int_{\R^d}\frac{\left|H_{jl}\left(x;\w\right)-H_{jl}\left(y;\w\right)\right|^2}{|x-y|^{d+2-\alpha}}dy\right)dx\\
&=\e_m^{-(2-\alpha)}\e_m^d\|h\|_\infty^2\int_{B\left(0,\frac{2R_0}{\e_m}\right)}
\left(\int_{\R^d}\frac{\left|\tilde H_{jl}(\tau_{x+z}\w)-\tilde H_{jl}(\tau_{x}\w)\right|^2}{|z|^{d+2-\alpha}}dz\right)dx.
\end{align*}
According to ergodic theorem we derive
\begin{equation}\label{t2-1-8}
\begin{split}
&\quad \lim_{m \to \infty}\left|B\left(0,\frac{2R_0}{\e_m}\right)\right|^{-1}\int_{B\left(0,\frac{2R_0}{\e_m}\right)}
\left(\int_{\R^d}\frac{\left|\tilde H_{jl}(\tau_{x+z}\w)-\tilde H_{jl}(\tau_{x}\w)\right|^2}{|z|^{d+2-\alpha}}dz\right)dx\\
&=\Ee\left[\int_{\R^d}\frac{\left|\tilde H_{jl}(\tau_z \w)-\tilde H_{jl}(\w)\right|^2}{|z|^{d+2-\alpha}}dz\right]
\le \|\tilde H_{jl}\|_{1-\frac{\alpha}{2}}^2,
\end{split}
\end{equation}
which implies immediately that
\begin{align*}
\int_{\R^d}\int_{\R^d}\frac{\left|H_{jl}\left(\frac{x}{\e_m};\w\right)-H_{jl}\left(\frac{y}{\e_m};\w\right)\right|^2|h(y)|^2}{|x-y|^{d+2-\alpha}}dxdy
%\le c_{17}(\w)\e_m^{-(2-\alpha)}\|\tilde H_{jl}\|_{1-\frac{\alpha}{2}}^2\le
\le c_{18}(\w)\e_m^{-(2-\alpha)},
\end{align*}
where we also use the fact $\|\tilde H_{jl}\|_{1-\frac{\alpha}{2}}^2\le\|\tilde H_{jl}\|_{1}^2<\infty$ and $c_{18}(\w)$
is a positive constant which may depend on $\w$.

Using ergodic theorem we have
\begin{align*}
\varlimsup_{m \to \infty}\int_{B(0,2R_0)}\frac{\left|H_{jl}\left(\frac{x}{\e_m};\w\right)\right|^2}{(1+|x|)^{d+2-\alpha}}dx
&\le c_{19}\varlimsup_{m \to \infty}\int_{B(0,2R_0)}\left|\tilde H_{jl}\left(\tau_{\frac{x}{\e_m}}\w\right)\right|^2dx
=c_{19}|B(0,2R_0)|\cdot\Ee[|\tilde H_{jl}|^2].
\end{align*}
And by direct computation we obtain
\begin{align*}
&\quad \int_{B(0,2R_0)^c}\frac{\left|H_{jl}\left(\frac{x}{\e_m};\w\right)\right|^2}{(1+|x|)^{d+2-\alpha}}dx
=\frac{1}{|B(0,1)|}\int_{B(0,2R_0)^c}\left(\int_{B(0,1)}\frac{\left|H_{jl}\left(\frac{x}{\e_m};\w\right)\right|^2}{(1+|x|)^{d+2-\alpha}}dy\right)dx\\
&\le c_{20}\int_{B(0,2R_0)^c}\int_{B(0,1)}\frac{\left|H_{jl}\left(\frac{x}{\e_m};\w\right)-H_{jl}\left(\frac{y}{\e_m};\w\right)\right|^2}{|x-y|^{d+2-\alpha}}dydx\\
&\quad +c_{20}\int_{B(0,1)}\left|H_{jl}\left(\frac{y}{\e_m};\w\right)\right|^2
\left(\int_{B(0,2R_0)^c}\frac{1}{(1+|x|)^{d+2-\alpha}}dx\right)dy\\
&\le c_{21}\e_m^{-(2-\alpha)}\e_m^d\int_{B\left(0,\frac{1}{\e_m}\right)}
\left(\int_{\R^d}\frac{\left|\tilde H_{jl}(\tau_{x}\w)-\tilde H_{jl}(\tau_{y}\w)\right|^2}{|x-y|^{d+2-\alpha}}dx\right)dy
+c_{21}\int_{B(0,1)}\left|\tilde H_{jl}\left(\tau_{\frac{y}{\e_m}}\w\right)\right|^2dy,
\end{align*}
where in the last step we have used the change of variable.
Therefore based on this estimate and applying the same arguments of \eqref{t2-1-8} we obtain
\begin{align*}
\int_{B(0,2R_0)^c}\frac{\left|H_{jl}\left(\frac{x}{\e_m};\w\right)\right|^2}{(1+|x|)^{d+2-\alpha}}dx
\le c_{22}(\w)\e_m^{-(2-\alpha)}.
\end{align*}
Hence combining all above estimates together yields that for every $m\ge 1$,
\begin{align*}
\int_{\R^d}|\xi|^{2-\alpha}\left|\F\left(H_{jl}^{\e_m}(\cdot;\w)h(\cdot)\right)(\xi)\right|^2d\xi
\le c_{23}(\w)\e_m^{-(2-\alpha)}.
\end{align*}
According to  this and \eqref{t1-2-9} together we derive
\begin{align*}
\e_m^{2-\alpha}\left|\int_{\R^d}\frac{\partial}{\partial x_l}\left( H_{jl}^{\e_m}\left(\cdot;\w\right)h(\cdot)\right)(x)
\frac{\partial g(x)}{\partial x_j}u_{\lambda,f}^{\e_m,\w}(x)dx\right|\le c_{24}(\w)\e_m^{1-\alpha/2}.
\end{align*}
This, along with \eqref{t1-2-7a} yields \eqref{t1-2-7} and we have finished the proof.

\end{proof}

\begin{thebibliography}{99}
\bibitem[CCKW1]{CCKW1} X. Chen, Z.-Q. Chen, T. Kumagai and J. Wang: Homogenization of symmetric stable-like processes in stationary ergodic media, \emph{SIAM J. Math. Anal.}, \textbf{53}, 2021, 2957--3001. %
\bibitem[JKO]{JKO} V.V. Jikov, S.M. Kozlov and O.A. Oleinik: \emph{Homogenization of Differential Operators and Integral Functionals}, Springer-Verlag, Berlin, 1994. %
\end{thebibliography}
\end{document}
